\documentclass[aspectratio=169,10pt]{beamer}
\usetheme[numbering=fraction,progressbar=frametitle]{metropolis}
\usepackage{fontspec}
\setsansfont{Latin Modern Sans}
\usepackage{amsmath,amssymb,mathtools,graphicx,booktabs}
\usepackage{polyglossia}
\setdefaultlanguage{english}
\setbeamertemplate{navigation symbols}{}
\setbeamerfont{frametitle}{size=\large}
\setbeamerfont{normal text}{size=\normalsize}
\setlength{\parskip}{.45em}
\hypersetup{pdfauthor={},pdfsubject={Lecture slides},pdfkeywords={},pdftitle={Matrix sequences and Markov chains}}
\providecommand{\N}{\mathbb{N}}
\providecommand{\Z}{\mathbb{Z}}
\providecommand{\D}{\mathbb{D}}
\providecommand{\Q}{\mathbb{Q}}
\providecommand{\R}{\mathbb{R}}
\providecommand{\C}{\mathbb{C}}
\providecommand{\vect}[1]{\overrightarrow{#1}}
\providecommand{\vi}{\vec{\imath}}
\providecommand{\vj}{\vec{\jmath}}
\providecommand{\vk}{\vec{k}}
\providecommand{\norm}[1]{\left\lVert #1\right\rVert}
\providecommand{\vc}[2]{\begin{pmatrix}#1\\#2\end{pmatrix}}
\providecommand{\deter}[2]{\det\left(#1,#2\right)}
\providecommand{\abs}[1]{\left\lvert #1\right\rvert}
\providecommand{\dd}{\,\mathrm{d}}
\providecommand{\e}{\mathrm{e}}
\providecommand{\Card}{\operatorname{Card}}
\providecommand{\E}{\operatorname{E}}
\providecommand{\V}{\operatorname{V}}
\providecommand{\degre}{^{\circ}}
\providecommand{\un}[1]{\,\mathrm{#1}}
\providecommand{\sci}[2]{#1\times10^{#2}}
\providecommand{\pgcd}{\operatorname{PGCD}}
\providecommand{\ppcm}{\operatorname{PPCM}}
\title{Matrix sequences and Markov chains}
\subtitle{Lecture and worked examples}
\date{}
\author{}
\begin{document}
\begin{frame}[plain]
\vfill
{\usebeamerfont{title}\inserttitle\par}
\medskip
{\usebeamerfont{subtitle}\insertsubtitle\par}
\medskip
{\usebeamercolor[fg]{progress bar}\rule{\textwidth}{.5pt}}
\vfill
\end{frame}
\begin{frame}{Objectives}
Study a sequence of column matrices defined by \(U_{n+1}=AU_n+B\); model a random evolution with two or three states by a probabilistic graph and a transition matrix; compute the distribution after \(n\) steps; find an invariant distribution; study convergence in the two-state case.
\end{frame}
\begin{frame}{Introductory activity}
In a simplified model (fictitious data), the weather on a given day is either fine (\(S\)) or rainy (\(P\)), and depends only on the weather the day before: after a fine day, it is fine with probability \(0.8\); after a rainy day, it is fine with probability \(0.4\). Day \(0\) is fine. Let \(s_n\) and \(p_n\) be the probabilities that day \(n\) is fine or rainy.
\end{frame}
\begin{frame}{Introductory activity (2)}
\begin{enumerate}
\setcounter{enumi}{0}\item Represent the situation by a diagram: two vertices \(S\) and \(P\), an arrow for each possible change, carrying its probability.

\item Using a tree, compute \(s_1\), \(p_1\), then \(s_2\), \(p_2\).

\item Justify that \(s_{n+1}=0.8s_n+0.4p_n\) and \(p_{n+1}=0.2s_n+0.6p_n\).

\item Let \(\pi_n=\begin{pmatrix}s_n&p_n\end{pmatrix}\) and \(T=\begin{pmatrix}0.8&0.2\\0.4&0.6\end{pmatrix}\). Check that \(\pi_{n+1}=\pi_nT\). Compute \(\pi_{10}\) with a calculator. Start again with a rainy day \(0\). What do you notice?

\end{enumerate}
\end{frame}
\begin{frame}{Introductory activity (3)}
\begin{enumerate}
\setcounter{enumi}{4}\item Find the distribution \(\pi=\begin{pmatrix}a&b\end{pmatrix}\), with \(a+b=1\), that satisfies \(\pi T=\pi\). Compare it with the observations of question 4.

\end{enumerate}
\end{frame}
\begin{frame}{Introductory activity (4)}
\emph{Observing is not proving: the properties conjectured here are established in the course.}
\end{frame}
\begin{frame}{Matrix sequences}
Let \(A\) be a square matrix of order \(p\), \(B\) a column matrix and \((U_n)\) a sequence of column matrices such that \(U_{n+1}=AU_n+B\) for every \(n\). If \(I_p-A\) is invertible, the matrix \(C=(I_p-A)^{-1}B\) satisfies \(C=AC+B\), and for every natural number \(n\):
\[U_n=A^n\left(U_0-C\right)+C.\]
\end{frame}
\begin{frame}{Matrix sequences}
\((I_p-A)C=B\) can be written \(C=AC+B\). Subtracting: \(U_{n+1}-C=A(U_n-C)\). The sequence \(V_n=U_n-C\) satisfies \(V_{n+1}=AV_n\), so \(V_n=A^nV_0\) (induction), and \(U_n=A^n(U_0-C)+C\).
\end{frame}
\begin{frame}{Matrix sequences}
If all the entries of \(A^n\) tend to \(0\), then \(U_n\) converges to \(C\) (entry by entry). This behaviour can be observed with a calculator; it is taken for granted when the statement says so.
\end{frame}
\begin{frame}{Markov chains}
A system is observed at times \(0,1,2,\dots\); at each time, it is in one of the states \(1,\dots,k\) (\(k=2\) or \(3\)). Let \(X_n\) be the state at time \(n\). We say that \((X_n)\) is a (homogeneous) \emph{Markov chain} if the probability of moving from state \(i\) to state \(j\) depends neither on \(n\) nor on the earlier states: for all \(i\), \(j\) and \(n\) such that \(P(X_n=i)\neq0\),
\[P_{(X_n=i)}(X_{n+1}=j)=t_{ij}.\]
The numbers \(t_{ij}\) are the \emph{transition probabilities}.
\end{frame}
\begin{frame}{Markov chains}
The chain is represented by a \emph{probabilistic graph}: a weighted directed graph whose vertices are the states, with an arc from \(i\) to \(j\), of weight \(t_{ij}\), when \(t_{ij}>0\). The \emph{transition matrix} is the square matrix \(T=(t_{ij})\). Its entries are nonnegative and the sum of each row is \(1\).
\end{frame}
\begin{frame}{Markov chains}
The \emph{distribution} at time \(n\) is the row matrix \(\pi_n=\begin{pmatrix}P(X_n=1)&\cdots&P(X_n=k)\end{pmatrix}\); \(\pi_0\) is the initial distribution.
\end{frame}
\begin{frame}{Markov chains}
For every natural number \(n\): \(\pi_{n+1}=\pi_nT\) and \(\pi_n=\pi_0T^n\).
\end{frame}
\begin{frame}{Markov chains}
By the law of total probability, with the partition \(\left(X_n=1\right),\dots,\left(X_n=k\right)\) (events of probability zero are left out, as they contribute \(0\)):
\[P(X_{n+1}=j)=\sum_{i=1}^{k}P(X_n=i)\,t_{ij},\]
which is entry \(j\) of \(\pi_nT\). The second formula follows by induction.
\end{frame}
\begin{frame}{Markov chains}
If \(P(X_0=i)=1\), then for every state \(j\), \(P(X_n=j)\) is the entry in row \(i\) and column \(j\) of \(T^n\). In other words, the probability of moving from state \(i\) to state \(j\) in \(n\) transitions is \(\left(T^n\right)_{ij}\).
\end{frame}
\begin{frame}{Markov chains}
The initial distribution is the row \(\pi_0\) with a \(1\) in position \(i\) and \(0\)s elsewhere. Then \(\pi_n=\pi_0T^n\) is row \(i\) of \(T^n\).
\end{frame}
\begin{frame}{Markov chains}
A distribution \(\pi\) (a row of nonnegative entries with sum \(1\)) is \emph{invariant} if \(\pi T=\pi\). If \(\pi_0=\pi\), then \(\pi_n=\pi\) for every \(n\).
\end{frame}
\begin{frame}{Markov chains}
Let \(T=\begin{pmatrix}1-a&a\\b&1-b\end{pmatrix}\) with \(0\le a\le1\), \(0\le b\le1\) and \(0<a+b<2\). The chain has a unique invariant distribution \(\pi=\begin{pmatrix}\dfrac b{a+b}&\dfrac a{a+b}\end{pmatrix}\), and whatever the initial distribution, \(\pi_n\) converges to \(\pi\).
\end{frame}
\begin{frame}{Markov chains}
Write \(\pi_n=\begin{pmatrix}x_n&1-x_n\end{pmatrix}\). Then \(x_{n+1}=(1-a)x_n+b(1-x_n)=(1-a-b)x_n+b\). This arithmetico-geometric sequence has fixed point \(\ell=\dfrac b{a+b}\), and \(x_n-\ell=(1-a-b)^n(x_0-\ell)\). Since \(-1<1-a-b<1\), \((1-a-b)^n\to0\) and \(x_n\to\ell\). An invariant distribution satisfies \(x=(1-a-b)x+b\), so \(x=\ell\): it is unique.
\end{frame}
\begin{frame}{Markov chains}
Solve the system formed by \(\pi T=\pi\) and the condition \(a+b+c=1\); one of the three equations of \(\pi T=\pi\) follows from the other two. Possible convergence of \(\pi_n\) is observed with a calculator.
\end{frame}
\begin{frame}{Markov chains}
The distribution is a row and we multiply on the right by \(T\): \(\pi_{n+1}=\pi_nT\), not \(T\pi_n\). It is the rows of \(T\), not its columns, that add up to \(1\).
\end{frame}
\end{document}
